\section{Discussion}\label{sec:discussion}

\paragraph{How the problem was parallelised.} There are three levels, each with its
own technique. \emph{Across the sequence}, token positions are divided among MPI
ranks: a rank owns the queries, keys, values and running state of its positions,
and only keys and values travel, by all-gather, ring or one-sided pull.
\emph{Inside a rank}, the (head, query tile) jobs are divided among OpenMP threads
with a static schedule; a job writes only its own rows of state, so threads need
no locks, and only the main thread calls MPI. On a GPU the same jobs become thread
blocks of 32 query rows. \emph{Inside a tile}, the pairs are vectorised: sixteen
single-precision lanes per instruction on the CPU, one warp of 32 threads per
query row on the GPU. Only the first level communicates; the other two share
memory and need no synchronisation beyond the end of a parallel loop or a kernel.

\paragraph{Where the time goes.} The measurements separate three costs that a
single efficiency number mixes.
\begin{itemize}
\item \textbf{The serial kernel.} Everything else is measured against it, so it
was tuned first. Most of the gain came from two decisions that are not about the
algorithm: storing every array on a cache-line boundary (a factor of
\n{align.speedup}) and letting the compiler use the full vector width
(\n{serial.w512.gain.pct}\%).
\item \textbf{The machine under load.} Inside one node, for shards of more than a
few hundred rows, almost all of the loss is already in the control. Hardware
counters split it into two factors: the clock, \n{clock.fratio.72} when a socket is
full, and the work per cycle, which follows how finely the sequence is cut
(\n{clock.wratio.36} at 1{,}820 rows per rank, \n{clock.wratio.72} at 910; contention
between the cores of a full node may contribute to it). Their
product, not communication, is the efficiency of the one-node runs with large
shards.
\item \textbf{Communication.} For a given configuration it is governed by the rows
a rank owns. With single-thread ranks inside a node a two-sided ring loses half its
speed at about \n{model.cstar.fit} rows per rank, and the fitted form puts the 5\%
mark near \n{model.c95.fit} rows. Most of that constant is explained by what the
exchange costs when all ranks of the node exchange at once (\n{ringx.cstar} rows
from a microbenchmark of exactly that) instead of what a message costs between two
ranks of an idle node (\n{model.cstar}); for shards below about 256 rows the
prediction from that measurement is still too pessimistic. Across nodes the overlapped ring stays close to
the same curve, the blocking ring and the all-gather lose more
(\n{strong.Ecomm.ring.big.288} and \n{strong.Ecomm.allgather.big.288} instead of
\n{strong.Ecomm.big.288} at 288 ranks for $N=131{,}072$), and the one-sided variants
become \n{strong.slow.rma_get.big.288} to \n{strong.slow.rma_get.small.288} times
slower than a run without communication. Faster ranks (threads, GPUs) move the
curve towards larger shards unless the transfer overlaps with the computation,
which it did for GPUs and, across nodes, did not for CPU ranks.
\end{itemize}

\paragraph{Which parts communicate most.} Per ring step a rank receives one shard
of keys and values, $(d+d_v)H_{kv}\,c$ numbers, and uses it for $c^2H$ query--key
pairs; nothing else is exchanged, and queries and state never move. The traffic of
a whole run is about $N(d+d_v)H_{kv}$ numbers per rank whatever the strategy, so the
strategies differ in \emph{when} and \emph{how} the bytes move, not in how many:
the all-gather moves them at once and stores all of them (\n{mem.ratio} times the
memory of the ring at 72 ranks and $N=\n{mem.N.h}$); the ring moves one shard per step; the one-sided
pull can skip shards that a mask makes useless, which halves the traffic for a
contiguous causal partition and saves nothing for the balanced one (a saving that
outweighed the imbalance only across nodes, where the pull is slow in any case). Grouped-query
attention is the one lever that reduces the bytes themselves, by the group factor
$G$.

\paragraph{Design choices that mattered.}
\begin{itemize}
\item \emph{One compute core, many ways to feed it.} Because every variant calls
the same kernel on the same blocks, two variants differ only in how the blocks are
moved, stored and synchronised. The no-communication control M0 gives that
difference a scale: $T_\text{M0}/T$ is the speed of a variant relative to a run that
exchanges nothing, and $T_S/(P\,T_\text{M0})$ is the efficiency of that run. The
control is not a perfect isolator (it holds the whole sequence and copies blocks
instead of receiving them), but the kernel time of both agreed within a few per
cent.
\item \emph{Online softmax as the interface.} A constant-size state per query that
can absorb key blocks in any order is what allows a ring, a pull in arbitrary
order and a GPU kernel to share one definition of correctness.
\item \emph{Partition by work, not by rows.} With a causal mask the zigzag
partition makes a ring \n{causal.zigzag.gain.overlap.288} to
\n{causal.zigzag.gain.overlap.72} times faster on the CPU, and
\n{gpu2.one.causal.gain} times with two GPUs, where a contiguous partition is less
unbalanced. For the
one-sided pull across nodes the contiguous partition was the faster of the two,
because it moves half the bytes.
\item \emph{Measure the measurement.} Two artefacts, a clock that needs seconds to
rise and arrays that the allocator places differently for every rank count, each
produced plausible but wrong scaling curves, including a superlinear speedup. Both
were found by asking why a number looked too good, and both are now controlled by
the protocol and demonstrated by their own experiments.
\end{itemize}

\paragraph{What I would choose.} For a sequence of $N$ tokens on $P$ ranks of this
machine: inside one node any of the strategies works, and the one-sided pull was
the cheapest when ranks own few rows; across nodes use the overlapped ring and
keep a rank above a few hundred rows, if necessary by using fewer ranks with more
threads; do not all-gather when memory per rank matters; with a causal mask use
the zigzag partition with a ring. The last two pieces of advice do not combine in
this implementation: with a causal mask, threads inside a rank are not balanced
(see the limits below), so a causal run should use one thread per rank. The
fastest configuration measured on four
nodes for $N=131{,}072$ was \n{hybrid.best.ranks.288} ranks of
\n{hybrid.best.threads.288} threads with the overlapped ring:
\n{hybrid.best.time.288}\,s, a speedup of \n{hybrid.speedup.288}. Two GPUs do the
same problem in \n{gpu2.one.T.host.overlap.N131072}\,s in one node and
\n{gpu2.two.T.host.overlap.N131072}\,s across two. These are choices for this
cluster and these shapes, not rules for other machines.

\paragraph{Limits of the evidence.} All results come from one cluster whose MPI
installation uses TCP between nodes. That a transfer does not advance while the
peer computes was shown for that installation with a two-rank microbenchmark and
is consistent with the application runs; it is not a property of MPI, and an
interconnect with remote memory access in hardware could reverse it. The
decomposition into $E_0$ and $E_\text{comm}$ is an identity around a control, not a
separation of causes; the split of $E_0$ into clock and work per cycle is measured,
but why the clock falls (power, temperature) is not. The fitted curve describes
single-thread ranks inside one node, and outside that setting rows per rank are
one variable among several. Every
configuration was measured three times in a row within one job, with one seed, so
differences of a few per cent are not resolved. The partition is static; the four
nodes used were equally fast, so its behaviour on unequal nodes was not tested.
The largest CPU run uses 288 cores and the largest GPU run two devices. The
threads of a rank take contiguous runs of its query tiles. That balances an
unmasked run, but with a causal mask the late rows of a rank carry more work than
the early ones, so the zigzag partition balances the ranks and not the threads
inside a rank; a development check on four cores showed most of the causal saving
lost when threads replaced ranks. No timing in this report combines a causal mask
with more than one thread per rank, and an interleaved or dynamic assignment of
tiles would be needed before doing so. The GPU
kernels were not profiled, GPU runs warm up with one repetition only, and what
limits the exchange between GPU nodes at small sizes was not found. Only the
forward pass of one layer with as many queries as keys was studied: not training,
and not decoding with few queries and a long key cache, where the balance of
computation and communication is different. The real-data conclusions rest on one
model, one text and three layers.

\paragraph{A design that was not explored.} The all-gather and the ring are the two
ends of a scale: every rank holds a copy of everything, or no rank holds more than
two shards. For all-pairs interactions it is known that the range in between is
useful: if every block is replicated on $r$ ranks, the number of messages falls
by $r^2$ and their volume by $r$, at $r$ times the memory~\cite{canbody}.
Attention with the sequence split by position is such an all-pairs problem, as the
comparison with the kernel matrix of Lecture~8 suggests, so a ring over groups of
ranks that each hold $r$ shards would trade memory for communication exactly
where \cref{fig:collapse} says communication hurts, at few rows per rank. I did not
implement it. Other things left for later are named in the limits above: a
profile and further tuning of the GPU kernels, more than two GPUs, the backward
pass, and a network that can move data without the peer's help.

\paragraph{What exact attention showed about the data.} In a real 24-layer model
at a context of 32{,}768 tokens the three layers examined behave in three ways:
the first is local, a middle layer has heads that put most of their probability on
the first token of the sequence, and the last is diffuse. For the windows tested, a
sliding window alone therefore changes the output of the middle layer completely,
keeping a few sink tokens repairs that layer, and even then a 1{,}024-token window
changes the output of the last layer by a third. These are changes of single
attention outputs for fixed inputs, not of what the model finally predicts. The
statistics came from a double-precision analysis of sampled queries; the
distributed code reproduced the outputs of the same layers (\cref{tab:v7}), which
is what one would run to obtain the exact output for every query at this length
and beyond.
